% Macro locali del problema; estratte da latex/SA_Book.tex.
\providecommand{\distanza}{}

\providecommand{\thetazero}{}

\providecommand{\phis}{}

\providecommand{\AreaIntPost}[1]{
\psset{unit=5ex}
\pspicture(-0.6,-2.5)(5.3,2.5)
\renewcommand{\distanza}{#1}
\renewcommand{\thetazero}{4.5 \distanza\space 2 \distanza\space 2 exp sub
sqrt atan 10 div sub}
\renewcommand{\phis}{4.5 \distanza\space 2 exp 2 mul 1 sub sqrt 1 atan
10 div sub}
\psplot[linewidth=2\pslinewidth]{\thetazero}{4.5}{10 x mul cos 10 x
mul sin add \distanza\space sub 2 exp 20 x mul sin 2 mul div neg 1 add
\distanza\space sub 60 mul}
\psplot[linewidth=2\pslinewidth]{0}{\thetazero}{\distanza\space 10 x
mul cos div \distanza\space sub 60 mul}
\psline[linewidth=0.7\pslinewidth]{->}(-0.3,0)(4.9,0)
\psline[linewidth=0.7\pslinewidth]{->}(0,-1)(0,2)
\psdots[dotstyle=|](!\thetazero\space 0)
\psdots[dotstyle=|](!\phis\space 0)
\psdots[dotstyle=|](4.5,0)
\uput[-90](!\thetazero\space 0){$\theta_{0}$}
\uput[-90](!\phis\space 0){$\theta_{s}$}
\uput[-90](4.5,0){$\pi/4$}
\rput(2,-2){$d=#1$}
\psframe[linecolor=cyan](-0.6,-1.5)(5.3,2.5)
\endpspicture
}

\providecommand{\AreaIntPre}[1]{
\psset{unit=5ex}
\pspicture(-0.6,-2.5)(5.3,5.5)
\renewcommand{\distanza}{#1}
\renewcommand{\thetazero}{4.5 \distanza\space 2 \distanza\space 2 exp sub
sqrt atan 10 div sub}
\psplot[linewidth=2\pslinewidth]{\thetazero}{4.5}{10 x mul cos 10 x
mul sin add \distanza\space sub 2 exp 20 x mul sin 2 mul div neg 1 add
\distanza\space sub 60 mul}
\psplot[linewidth=2\pslinewidth]{0}{\thetazero}{\distanza\space 10 x
mul cos div \distanza\space sub 60 mul}
\psline[linewidth=0.7\pslinewidth]{->}(-0.3,0)(4.9,0)
\psline[linewidth=0.7\pslinewidth]{->}(0,-1)(0,5)
\psdots[dotstyle=|](!\thetazero\space 0)
\psdots[dotstyle=|](4.5,0)
\uput[-90](!\thetazero\space 0){$\theta_{0}$}
\uput[-90](4.5,0){$\pi/4$}
\rput(2,-2){$d=#1$}
\psframe[linecolor=cyan](-0.6,-1.5)(5.3,5.5)
\endpspicture
}

\renewcommand{\Anno}{2010}

\renewcommand{\NumeroProblema}{4}

\renewcommand{\Testo}{
Definiamo \emph{striscia} di ampiezza $d$ la parte di piano compresa tra 2 rette parallele poste a distanza $d$ l'una dall'altra. Consideriamo ora una striscia di ampiezza $d<1$ ed un quadrato di lato 1 con 2 lati paralleli alle rette che delimitano la striscia, ed il centro equidistante dalle rette stesse.

Successivamente ruotiamo il quadrato in senso orario di un angolo $\theta$ (rispetto al centro).

\begin{enumerate}
	\item Determinare, in funzione di $d$, il più piccolo valore di $\theta$ per cui 2 vertici del quadrato vengono a trovarsi sulle rette che delimitano la striscia.

	\item Detto $\theta_{0}$ il valore trovato al punto precedente, determinare (in funzione di $d$ e $\theta$) l'area dell'intersezione tra il quadrato e la striscia per valori di $\theta$ tali che $0\leq\theta\leq\theta_{0}$.

	\item Calcolare (in funzione di $d$) l'area dell'intersezione nel caso $\theta=45^{\circ}$.

	\item Detto $\theta_{0}$ il valore trovato al punto (a), determinare (in funzione di $d$ e $\theta$) l'area dell'intersezione tra il quadrato e la striscia per valori di $\theta$ tali che $\theta_{0}\leq\theta\leq 45^{\circ}$.

\end{enumerate}
}

\renewcommand{\Risposta}{
\begin{enumerate}
	\item[(a)]  Ci sono tantissimi modi (equivalenti) di esprimere la risposta. Ne segnaliamo \emph{alcuni}:
	$$
	\begin{array}{c@{\hspace{3em}}c@{\hspace{3em}}c}
		\displaystyle{\frac{\pi}{4}-\arcsin\frac{d}{\sqrt{2}},} &
		\displaystyle{-\frac{\pi}{4}+\arccos\frac{d}{\sqrt{2}},} &
		\displaystyle{\frac{\pi}{4}-\arctan\frac{d}{\sqrt{2-d^{2}}},} \\
		\noalign{\vspace{2ex}}
		\displaystyle{\frac{1}{2}\arcsin(1-d^{2}),} &
		\displaystyle{\arcsin\frac{\sqrt{2-d^{2}}-d}{2},} &
		\displaystyle{\arccos\frac{\sqrt{2-d^{2}}+d}{2}.}
	\end{array}$$

	\item[(b)]  L'area dell'intersezione \`{e} $d/\cos\theta$.

	\item[(c)]  L'area dell'intersezione \`{e}
	$$1-\frac{\left(\sqrt{2}-d\right)^{2}}{2}=
	\sqrt{2}\,d-\frac{d^{2}}{2}.$$

	\item[(d)]  L'area dell'intersezione \`{e}
	$$1-
\frac{\left(\cos\theta+\sin\theta-d\right)^{2}}{2\sin(2\theta)}.$$
\end{enumerate}
}

\renewcommand{\Aiutino}{
\begin{enumerate}
	\renewcommand{\labelenumi}{(\alph{enumi})} \item[(a)] Ragionare su uno
	dei tanti triangoli rettangoli presenti nella figura, ad esempio
	quello che ha come vertici il centro del quadrato, uno dei vertici
	del quadrato che stanno su una delle rette che delimitano la
	striscia, e la proiezione del centro sulla retta stessa. Quanto 
	misurano gli angoli di questo triangolo?

	\item[(b)] L'intersezione \`{e} un parallelogrammo che ha altezza $d$.
	Non resta che esprimere la lunghezza della base in funzione di 
	$\theta$. 

	\item[(c)] Conviene calcolare l'area dei due triangolini (uguali) che
	spuntano dalla striscia. La loro altezza si calcola facilmente 
	pensando alla diagonale del quadrato. Quanto misurano gli angoli 
	dei due triangolini?

	\item[(d)] Come nel caso precedente, conviene calcolare l'area dei due
	triangolini (uguali) che spuntano dalla striscia.  Quanto 
	misurano i loro angoli? E la loro altezza? Quale angolo forma la 
	diagonale del quadrato (quella che congiunge i due vertici 
	esterni) con una retta perpendicolare alla striscia?
\end{enumerate}
}

\renewcommand{\SoluzioneUfficiale}{\TitoloSoluzione \DomandaA

Con riferimento alla Figura~\ref{fig:quadrato-1}, consideriamo il
triangolo rettangolo $H\!AO$.  Sappiamo che il cateto $OH$ misura
$d/2$ e l'ipotenusa $OA$ misura $\sqrt{2}/2$ (mezza diagonale del
quadrato).  Inoltre l'angolo che il lato $DA$ forma con ciascuna delle
rette che delimitano la striscia \`{e} uguale all'angolo di rotazione
$\theta_{0}$, dunque
$$\angle OAH=\angle OAD-\angle HAD=45^{\circ}-\theta_{0}.$$

Per le note relazioni trigonometriche nei triangoli rettangoli si 
avr\`{a} pertanto che 
$$OA\cdot\sin(\angle OAH)=OH,$$ 
quindi 
\begin{equation}
	\frac{\sqrt{2}}{2}\cdot\sin\left(45^{\circ}-\theta_{0}\right)=
	\frac{d}{2},
	\label{eqn:theta-0}
\end{equation}
da cui si ricava facilmente che 
$$\theta_{0}=45^{\circ}-\arcsin\frac{d}{\sqrt{2}}.$$

\begin{figure}[hbt]
	\centering
	\psset{unit=3.5ex}
	\pspicture(-8,-5.5)(8,5.5)
	\newrgbcolor{azzurrino}{0.6 0.9 1}
	\newcommand{\angolo}{20}
	\pstGeonode[PointName=none,PointSymbol=none](0,0){O}(4,4){A'}(4,-4){B'}(-4,-4){C'}(-4,4){D'}
	\pstGeonode[PointName=none,PointSymbol=none](8,4){A1}(8,-4){B1}(-8,-4){C1}(-8,4){D1}
	\pstRotation[PosAngle=45,RotAngle=-\angolo]{O}{A'}[A]
	\pstRotation[RotAngle=-\angolo]{O}{B'}[B]
	\pstRotation[PosAngle=-90,RotAngle=-\angolo]{O}{C'}[C]
	\pstRotation[PosAngle=135,RotAngle=-\angolo]{O}{D'}[D]
	\pstProjection[PointName=none,PointSymbol=none]{A1}{B1}{A}[A2]
	\pstProjection[PointName=none,PointSymbol=none]{C1}{D1}{C}[C2]
	\psframe*[linecolor=azzurrino](C2)(A2)
	\pstProjection{A}{A2}{O}[H]
	\psline(A)(B)(C)(D)(A)
	\psline(O)(H)(A)(O)
	\psdots(O)(A)(C)
	\pstMarkAngle[MarkAngleRadius=1,LabelSep=2.2]{D}{A}{H}{$\theta_{0}$}
	\pstRightAngle[linewidth=0.5\pslinewidth,RightAngleSize=.5]{O}{H}{A}
	\uput[225](0,0){$O$}
	\endpspicture
	\caption{Caso in cui due vertici del quadrato stanno sulle rette 
	che delimitano la striscia}
	\label{fig:quadrato-1}
\end{figure}
\begin{center}
\end{center}


Si noti che, come ci si poteva facilmente aspettare, $\theta_{0}$ 
\`{e} sempre positivo e minore di $45^{\circ}$. Inoltre
\begin{itemize}
	\item  quando $d\to 1$ si ha che $\theta_{0}\to 0$, 

	\item quando $d\to 0$ (cio\`{e} la striscia diventa
	sostanzialmente una retta) si ha che $\theta_{0}\to 45^{\circ}$.
\end{itemize}

\DomandaB

Siamo nel caso in cui il punto $A$ si trova all'esterno della striscia.
Con riferimento alla Figura~\ref{fig:quadrato-2}, l'intersezione tra
il quadrato e la striscia \`{e} il parallelogrammo $PQRS$, la cui area
\`{e} $PQ\cdot d$.  Per calcolare $PQ$ osserviamo che \`{e} uguale
alla lunghezza del segmento $AT$, dove la retta $AT$ \`{e} costruita
parallelamente alla retta $PQ$.

\begin{figure}[htbp]
	\centering
	\psset{unit=3ex}
	\pspicture(-8,-5)(8,5.5)
	\newrgbcolor{azzurrino}{0.6 0.9 1}
	\newcommand{\angolo}{15}
	\pstGeonode[PointName=none,PointSymbol=none](0,0){O}(4,4){A'}(4,-4){B'}(-4,-4){C'}(-4,4){D'}
	\pstGeonode[PointName=none,PointSymbol=none](8,2){A1}(8,-2){B1}(-8,-2){C1}(-8,2){D1}
	\pstRotation[PosAngle=45,RotAngle=-\angolo]{O}{A'}[A]
	\pstRotation[RotAngle=-\angolo]{O}{B'}[B]
	\pstRotation[PosAngle=-90,RotAngle=-\angolo]{O}{C'}[C]
	\pstRotation[PosAngle=135,RotAngle=-\angolo]{O}{D'}[D]
	\psframe*[linecolor=azzurrino](C1)(A1)
	\pstInterLL[PosAngle=30]{A}{B}{D1}{A1}{P}
	\pstInterLL[PosAngle=140]{C}{D}{D1}{A1}{Q}
	\pstInterLL[PosAngle=-30]{A}{B}{C1}{B1}{S}
	\pstInterLL[PosAngle=-150]{C}{D}{C1}{B1}{R}
	\pstTranslation[PosAngle=140]{P}{Q}{A}[T]
	\pspolygon*[linecolor=cyan](P)(Q)(R)(S)
	\psline(A)(B)(C)(D)(A)
	\psline(T)(A)
	\psline(Q)(P)
	\psline(R)(S)
	\pstMarkAngle[MarkAngleRadius=1.3,LabelSep=2.4]{D}{A}{T}{$\theta$}
	\endpspicture
	\caption{caso in cui $0<\theta<\theta_{0}$}
	\label{fig:quadrato-2}
\end{figure}

Applicando le note relazioni trigonometriche nel triangolo rettangolo 
$DTA$, avremo pertanto che
$$AT\cdot\cos\theta=AD=1,$$
da cui
$$PQ=AT=\frac{1}{\cos\theta},$$
ed infine
$$\mbox{Area intersezione}=PQ\cdot d=\frac{d}{\cos\theta}.$$

Si noti che tale area \emph{aumenta} al crescere di $\theta$ 
(perch\'{e} il numeratore \`{e} fisso ed il denominatore diminuisce).

\DomandaC


Per calcolare l'area dell'intersezione, conviene calcolare l'area dei
due triangoli esclusi dall'intersezione stessa.  Con riferimento alla
Figura~\ref{fig:quadrato-3}, nel fare questo possiamo limitarci al
triangolo $DQP$, in quanto l'altro \`{e} uguale per simmetria.

Il triangolo $DQP$ \`{e} rettangolo isoscele, in quanto gli angoli in 
$Q$ e $P$ misurano $45^{\circ}$. La sua altezza $DH$ \`{e} uguale a 
$(\sqrt{2}-d)/2$, cio\`{e} met\`{a} della diagonale del quadrato meno 
met\`{a} dell'ampiezza della striscia. Essendo la base il doppio 
dell'altezza, l'area sar\`{a} pertanto
$$\mbox{Area}(DQP)=\frac{\left(\sqrt{2}-d\right)^{2}}{4}.$$

\begin{figure}[htbp]
	\centering
	\psset{unit=3ex}
	\pspicture(-8,-7)(8,7)
	\newrgbcolor{azzurrino}{0.6 0.9 1}
	\newcommand{\angolo}{45}
	\pstGeonode[PointName=none,PointSymbol=none](8,3.2){A1}(8,-3.2){B1}(-8,-3.2){C1}(-8,3.2){D1}
	\psframe*[linecolor=azzurrino](C1)(A1)
	\pstGeonode[PointName=none,PointSymbol=none](0,0){O}(4,4){A'}(4,-4){B'}(-4,-4){C'}(-4,4){D'}
	\pstRotation[PosAngle=0,RotAngle=-\angolo]{O}{A'}[A]
	\pstRotation[PosAngle=-90,RotAngle=-\angolo]{O}{B'}[B]
	\pstRotation[PosAngle=180,RotAngle=-\angolo]{O}{C'}[C]
	\pstRotation[PosAngle=90,RotAngle=-\angolo]{O}{D'}[D]
	\pstInterLL[PosAngle=45]{A}{D}{D1}{A1}{P}
	\pstInterLL[PosAngle=135]{C}{D}{D1}{A1}{Q}
	\psline(A)(B)(C)(D)(A)
	\psline(Q)(P)
	\psline(D)(B)
	\pstProjection[PosAngle=-45]{Q}{P}{D}[H]
	\pstMarkAngle[MarkAngleRadius=0.7,LabelSep=1.5]{D}{P}{Q}{$45^{\circ}$}
	\pstMarkAngle[MarkAngleRadius=0.7,LabelSep=1.5]{P}{Q}{D}{$45^{\circ}$}
	\endpspicture
	\caption{caso $\theta=45^{\circ}$}
	\label{fig:quadrato-3}
\end{figure}

Per differenza si ottiene quindi che
\begin{equation}
	\mbox{Area intersezione}=1-\frac{\left(\sqrt{2}-d\right)^{2}}{2}.
	\label{intersezione-45}
\end{equation}

\DomandaD

Siamo nel caso in cui il punto $A$ cade all'interno della striscia.
Con riferimento alla Figura~\ref{fig:quadrato-4} il punto fondamentale
\`{e}, come nel caso precedente, trovare l'area del triangolo $DQP$.
Si tratta di un triangolo rettangolo di cui conosciamo gli angoli in
$P$ (cio\`{e} $\theta$) ed in $Q$ (cio\`{e} $90^{\circ}-\theta$).  Per
concludere dobbiamo calcolare l'altezza $DH$.

\begin{figure}[htbp]
	\centering
	\psset{unit=3ex}
	\pspicture(-8,-7)(8,7)
	\newrgbcolor{azzurrino}{0.6 0.9 1}
	\newcommand{\angolo}{30}
	\pstGeonode[PointName=none,PointSymbol=none](8,3.2){A1}(8,-3.2){B1}(-8,-3.2){C1}(-8,3.2){D1}
	\psframe*[linecolor=azzurrino](C1)(A1)
	\pstGeonode[PointName=none,PointSymbol=none](0,0){O}(4,4){A'}(4,-4){B'}(-4,-4){C'}(-4,4){D'}
	\pstRotation[PosAngle=0,RotAngle=-\angolo]{O}{A'}[A]
	\pstRotation[PosAngle=-90,RotAngle=-\angolo]{O}{B'}[B]
	\pstRotation[PosAngle=180,RotAngle=-\angolo]{O}{C'}[C]
	\pstRotation[PosAngle=90,RotAngle=-\angolo]{O}{D'}[D]
	\pstInterLL[PosAngle=45]{A}{D}{D1}{A1}{P}
	\pstInterLL[PosAngle=135]{C}{D}{D1}{A1}{Q}
	\pstProjection[PointName=none,PointSymbol=none]{C1}{D1}{B}[B2]
	\pstProjection[PointName=none,PointSymbol=none]{A1}{B1}{B}[B3]
	\pstProjection{B}{B2}{D}[K]
	\psline(A)(B)(C)(D)(A)
	\psline(Q)(P)
	\psline(D)(B)
	\psline(D)(K)
	\psline[linewidth=0.7\pslinewidth](B2)(B3)
	\pstProjection[PosAngle=-135]{Q}{P}{D}[H]
	\pstMarkAngle[MarkAngleRadius=0.7,LabelSep=1.3]{D}{P}{Q}{$\theta$}
	\pstMarkAngle[MarkAngleRadius=0.7,LabelSep=1.2]{C}{D}{K}{$\theta$}
	\pstRightAngle[linewidth=0.5\pslinewidth,RightAngleSize=.5]{B}{K}{D}
	\endpspicture
	\caption{caso $\theta_{0}<\theta<45^{\circ}$}
	\label{fig:quadrato-4}
\end{figure}

Tracciamo la parallela per $B$ alle rette che delimitano la striscia,
ed indichiamo con $K$ la proiezione di $D$ su tale parallela.  \`{E}
chiaro per ragioni di simmetria che $DH=(DK-d)/2$.  Per calcolare
$DK$, osserviamo che nel triangolo rettangolo $KBD$ l'ipotenusa $DB$
\`{e} $\sqrt{2}$ (la diagonale del quadrato) e
$$\angle KDB=\angle CDB-\angle CDK=45^{\circ}-\theta.$$

Ne segue che
$$DH=\frac{1}{2}(DK-d)=
\frac{1}{2}\left(\sqrt{2}\cos(45^{\circ}-\theta)-d\right),$$
da cui, applicando la formula di sottrazione del coseno, si ha che
$$DH=\frac{\cos\theta+\sin\theta-d}{2}.$$

Applicando note formule ai triangoli rettangoli $DQP$, $DHQ$ e $DHP$ 
avremo pertanto che
$$\mbox{Area}(DQP)=\frac{1}{2}DQ\cdot DP=
\frac{1}{2}\frac{DH}{\cos\theta}\cdot\frac{DH}{\sin\theta}=
\frac{DH^{2}}{\sin(2\theta)}=
\frac{\left(\cos\theta+\sin\theta-d\right)^{2}}{4\sin(2\theta)}.$$

Per differenza si ottiene quindi che
$$\mbox{Area intersezione}=1-
\frac{\left(\cos\theta+\sin\theta-d\right)^{2}}{2\sin(2\theta)}.$$

\`{E} relativamente semplice verificare che questa formula
\begin{itemize}
	\item  coincide con quella trovata al punto~(c) quando 
	$\theta=45^{\circ}$,

	\item  coincide con quella trovata al punto~(b) quando 
	$\theta=\theta_{0}$, dove $\theta_{0}$ \`{e} l'angolo trovato al 
	punto~(a), cio\`{e} l'angolo per cui 
	$$d=\sqrt{2}\sin\left(45^{\circ}-\theta\right)=
	\cos\theta-\sin\theta.$$
\end{itemize}
}

\renewcommand{\SoluzioniAlternative}{\TitoloSoluzioneAlternativaA
Il punto (a) pu\`{o} essere risolto in modo del tutto equivalente 
esaminando uno qualunque dei triangoli rettangoli presenti in figura. 
Ci si riduce ad un'equazione simile alla (\ref{eqn:theta-0}), 
equazione che si pu\`{o} sempre portare nella forma
$$\cos\theta-\sin\theta=d.$$

Da qui si pu\`{o} mettere a sistema con la relazione trigonometrica 
fondamentale $\cos^{2}\theta+\sin^{2}\theta=1$, e quindi ricavare 
$\cos\theta$ e/o $\sin\theta$ e/o $\tan\theta$ risolvendo il sistema. 
Si ottengono cos\`{i} molte delle scritture equivalenti di 
$\theta_{0}$ indicate nel capitolo ``Risposte''. Verificarne 
l'equaivalenza \`{e} un utile esercizio di trigonometria.

Un approccio alternativo \`{e} basato sulla 
Figura~\ref{fig:quadrato-5}. Applicando il teorema di Carnot nel 
triangolo $OAA'$ si ha che
\begin{equation}
	[AA']^{2}=OA^{2}+[OA']^{2}-2\,OA\cdot OA'\cdot\cos\theta.
	\label{eqn:quadrato-carnot}
\end{equation}

\begin{figure}[htb]
	\centering
	\psset{unit=2.5ex}
	\pspicture(-8,-7.5)(8,7.5)
	\newrgbcolor{azzurrino}{0.6 0.9 1}
	\psframe*[linecolor=azzurrino](-8,-3)(8,3)
	\pstGeonode[PosAngle={45,135,-135,-45,-90}](5,5){A}(-5,5){B}(-5,-5){C}(5,-5){D}(0,0){O}
	\psline(A)(B)(C)(D)(A)
	\pstCircleOA[linestyle=dashed,linecolor=blue]{O}{A}
	\pstGeonode[PosAngle=60](!41 sqrt 3){A'}
	\pstRotation[PosAngle=110,RotAngle=90]{O}{A'}[B']
	\pstRotation[PosAngle=240,RotAngle=180]{O}{A'}[C']
	\pstRotation[PosAngle=-80,RotAngle=270]{O}{A'}[D']
	\psline[linecolor=magenta](A')(B')(C')(D')(A')
	\psline[linewidth=1.3\pslinewidth](O)(A)(A')(O)
	\pstMarkAngle[MarkAngleRadius=1.5,LabelSep=2.2]{A'}{O}{A}{$\theta$}
	\endpspicture
	\caption{Approccio mediante teorema di Carnot}
	\label{fig:quadrato-5}
\end{figure}

Ora $OA$ e $OA'$ sono mezza diagonale del quadrato. Per calcolare 
$AA'$ possiamo procedere per via analitica. Ponendo coordinate 
cartesiane in modo che sia $O=(0,0)$ e $A=(1/2,1/2)$, avremo che $A'$ 
si ottiene intersecando la retta di equazione $y=d/2$ (il bordo 
superiore della striscia) con la circonferenza di equazione 
$x^{2}+y^{2}=1/2$ (la circonferenza su cui si muovono i vertici del 
quadrato durante la rotazione). Risolvendo il sistema si ricava che
$$A'=\left(\frac{\sqrt{2-d^{2}}}{2},\frac{d}{2}\right).$$

\`{E} ora immediato calcolare la distanza $AA'$ e quindi ricavare 
$\cos\theta$ dalla (\ref{eqn:quadrato-carnot}).
}

\renewcommand{\Osservazioni}{\TitoloOsservazione

\`{E} interessante chiedersi per quale valore di $\theta$ l'area
dell'intersezione risulta massima.  Grazie alle formule ricavate nella
soluzione, questo \`{e} equivalente a calcolare il massimo della
funzione 
$$\mbox{Area}(\theta)=\left\{
\begin{array}{ll}
	\displaystyle{\frac{d}{\cos\theta}} & \mbox{se
	}0\leq\theta\leq\theta_{0}, \\
	\noalign{\vspace{2ex}}
	\displaystyle{1-\frac{\left(\cos\theta+
	\sin\theta-d\right)^{2}}{2\sin(2\theta)}}\mbox{\quad}
	& \mbox{se }\theta_{0}\leq\theta\leq\displaystyle{\frac{\pi}{4}},
\end{array}
\right.$$
al variare di $\theta$ in $[0,\pi/4]$.

La risposta, forse non cos\`{i} intuitiva, \`{e} che
\begin{itemize}
	\item  se $0<d\leq\sqrt{2}/2$ l'area massima si ottiene per 
	$\theta=\pi/4$, ed \`{e} data dalla formula 
	(\ref{intersezione-45}),

	\item se $\sqrt{2}/2<d\leq 1$ l'area massima \`{e} $(d^{2}+1)/2$ e
	si ottiene per 
	$$\theta=\frac{\pi}{4}-\arccos\frac{1}{d\sqrt{2}}.$$
\end{itemize}

Vediamo un accenno a come si dimostra questo fatto. 

Come abbiamo gi\`{a} osservato, per $\theta\in[0,\theta_{0}]$ l'area
dell'intersezione \`{e} una funzione crescente di $\theta$.  Ci basta
quindi massimizzare l'espressione valida in $[\theta_{0},\pi/4]$.

Con il cambio di variabili $\phi=\pi/4-\theta$, e qualche formula 
trigonometrica, ci si riduce a massimizzare la funzione
$$f(\phi)=\mbox{Area}\left(\frac{\pi}{4}-\phi\right)=
1-\frac{\left(\sqrt{2}\cos\phi-d\right)^{2}}{2\cos(2\phi)}$$
al variare di 
\begin{equation}
	\phi\in\left[0,\frac{\pi}{4}-\theta_{0}\right].
	\label{intervallo-phi}
\end{equation}

Per far questo utilizziamo gli strumenti del calcolo differenziale.
Con qualche ulteriore calcolo si ottiene che
$$f'(\phi)=\sin\phi\left(\sqrt{2}\cos\phi-d\right)
\left(2d\cos\phi-\sqrt{2}\right).$$

Esaminiamo il segno di tale espressione in tutto l'intervallo 
$(0,\pi/4)$. In tale intervallo si ha chiaramente che $\sin\phi>0$ e
$$\sqrt{2}\cos\phi-d>\sqrt{2}\cdot\frac{1}{\sqrt{2}}-1\geq 0.$$

Il segno di $f'(\phi)$ dipende quindi unicamente dal segno 
dell'ultimo fattore. Consideriamo ora due casi.
\begin{itemize}
	\item  Se $0<d\leq\sqrt{2}/2$ si ha che
	$$2d\cos\phi-\sqrt{2}<2\cdot\frac{\sqrt{2}}{2}\cdot 
	1-\sqrt{2}\leq 0.$$
	
	Ne segue che $f'(\phi)<0$ in tutto l'intervallo $(0,\pi/4)$, 
	dunque $f(\phi)$ \`{e} strettamente decrescente nell'intervallo 
	stesso. In particolare il massimo viene assunto per $\phi=0$ e di 
	conseguenza, tornando al problema originario, si ha che
	$$\max\left\{\mbox{Area}(\theta):
	0\leq\theta\leq\frac{\pi}{4}\right\}=
	\mbox{Area}\left(\frac{\pi}{4}\right)=
	1-\frac{\left(\sqrt{2}-d\right)^{2}}{2}.$$

	\item  Se $\sqrt{2}/2<d\leq 1$, allora il terzo fattore cambia 
	segno. In particolare, posto
	$$\phi_{s}=\arccos\frac{1}{d\sqrt{2}},$$
	si ha che il terzo fattore \`{e} positivo in $(0,\phi_{s})$ e
	negativo in $(\phi_{s},\pi/4)$.  Poich\'{e} lo stesso discorso
	vale per il segno di $f'(\phi)$, ne segue che $f(\phi)$, vista
	come funzione in $(0,\pi/4)$, ha massimo per $\phi=\phi_{s}$ (che
	si verifica facilmente appartenere a tale intervallo).  Resta
	per\`{o} da stabilire se $\phi_{s}$ appartiene all'intervallo
	(\ref{intervallo-phi}).  Data l'espressione di $\theta_{0}$,
	ci\`{o} \`{e} equivalente a chiedersi se
	$$\arccos\frac{1}{d\sqrt{2}}\leq\arcsin\frac{d}{\sqrt{2}}.$$
	
	Sfruttando la nota relazione $\arccos x=\arcsin\sqrt{1-x^{2}}$ ci si 
	riduce a
	$$\arcsin\frac{\sqrt{2d^{2}-1}}{d\sqrt{2}}\leq
	\arcsin\frac{d}{\sqrt{2}}.$$
	
	Eliminando l'arcoseno, che \`{e} una funzione crescente, si trova 
	una disequazione irrazionale che risulta soddisfatta per ogni 
	valore ammissibile di $d$. Ne segue che in tal caso il massimo di 
	$f(\phi)$, al variare di $\phi$ nell'intervallo 
	(\ref{intervallo-phi}), \`{e} sempre raggiunto in $\phi_{s}$, e 
	con semplici conti si vede che vale $(d^{2}+1)/2$. 
	
	Tornando al problema originario abbiamo quindi che in tal caso
	$$\max\left\{\mbox{Area}(\theta):
	0\leq\theta\leq\frac{\pi}{4}\right\}=
	\mbox{Area}\left(\frac{\pi}{4}-\phi_{s}\right)=
	\frac{d^{2}+1}{2}.$$
\end{itemize}


\begin{figure}[htbp]
	\hfill
	\AreaIntPre{0.5}
	\hfill
	\AreaIntPre{0.6}
	\hfill
	\AreaIntPre{0.71}
	\hfill\mbox{}

	\hfill
	\AreaIntPost{0.73}
	\hfill
	\AreaIntPost{0.82}
	\hfill
	\AreaIntPost{0.835}
	\hfill\mbox{}\vspace{-4ex}
	\caption{andamento del grafico della funzione
	$\mbox{Area}(\theta)-d$ per diversi valori del parametro $d$.
	Sono indicati anche $\theta_{0}$ ed il punto di massimo $\theta_{s}$
	(quando \`{e} interno all'intervallo)}
	\label{fig:grafici-area}
\end{figure}

I grafici della Figura~\ref{fig:grafici-area} mostrano il cambio di
comportamento qualitativo (dall'avere il massimo in un estremo
all'avercelo in un punto interno) dell'area dell'intersezione non
appena $d$ attraversa il valore critico $\sqrt{2}/2$.  Per poter
meglio confrontare i grafici abbiamo sottratto $d$ all'area
dell'intersezione (il che non altera la crescenza o decrescenza del 
grafico).

Concludiamo con alcune ulteriori osservazioni.
\begin{itemize}
	\item Finora abbiamo determinato il massimo dell'area
	dell'intersezione con il vincolo che il centro del quadrato fosse
	equidistante dalle due rette che delimitano la striscia.  Cosa
	accade se lasciamo libera la posizione del centro?
	Sostanzialmente non cambia nulla.  Si pu\`{o} infatti dimostrare
	abbastanza semplicemente che, a parit\`{a} di angolo di rotazione,
	una delle posizioni (o la posizione) che realizza l'area massima
	\`{e} proprio quella in cui il centro \`{e} equidistante dalle due
	rette.  Ci\`{o} non dice che sia unica: ad esempio, per valori di
	$\theta$ minori di $\theta_{0}$ (come nella
	Figura~\ref{fig:quadrato-2}), una piccola traslazione
	perpendicolare alla striscia lascia invariata l'area
	dell'intersezione.

	\item Considerando le posizioni in cui il centro \`{e}
	equidistante dalle rette che delimitano la striscia, per quale
	angolo di rotazione l'area dell'intersezione \`{e} \emph{minima}?
	Poich\'{e} abbiamo visto che la funzione $\mbox{Area}(\theta)$
	\`{e} sempre crescente, oppure prima crescente e poi decrescente,
	il suo valore minimo si realizza necessariamente in uno degli
	estremi dell'intervallo.  Per stabilire quale, basta confrontare
	il valore in 0, e cio\`{e} $d$, con il valore in $\pi/4$ dato
	dalla (\ref{intersezione-45}).  Risolvendo la semplice
	disequazione risultante, si trova che il minimo si realizza per
	$\theta=0$ quando $0<d<2(\sqrt{2}-1)$ e per $\theta=\pi/4$ quando
	$2(\sqrt{2}-1)<d\leq 1$.

	\item Dai grafici si intuisce che la funzione
	$\mbox{Area}(\theta)$ \`{e} non solo continua in
	$\theta=\theta_{0}$ (cosa che abbiamo gi\`{a} verificato), ma ha
	anche derivata continua nello stesso punto.  Un po' di calcoli
	elementari (che per brevit\`{a} non riportiamo) confermano questa
	intuizione, stabilendo anche che
	$$\mbox{Area}'(\theta_{0})=\tan\theta_{0}-\tan^{2}\theta_{0}.$$ 
\end{itemize} 


}

\renewcommand{\Commenti}{
}